\section{Predictions, measured, and an honest ledger}
\label{sec:ledger}

Table~\ref{tab:predictions} lists every claim in this paper that has a
number behind it, the domain it belongs to, the measured value, and the
script that owns it. Each script ends with \emph{gates}: conditions its
own result must meet, written down before the result was known. The table
is generated from the scripts' data files by the same program that writes
every number in the text, and that program refuses to run if any gate
fails; so a claim cannot outlive its evidence, and the numbers in the prose
cannot drift from the numbers in the data.

\begin{table}[t]
  \centering\small
  % GENERATED by paper/tools/numbers.mjs -- do not edit by hand.
\begin{tabular}{@{}p{0.40\linewidth}lp{0.22\linewidth}p{0.17\linewidth}@{}}
\toprule
Claim & Domain & Measured & Script \\
\midrule
A linear observer sees no beat in an additive superposition & any & $8\times 10^{-17}$ & \texttt{observer.mjs} \\
A squaring front end mints it at half the product of the harmonics & any & 0.0403 vs 0.0403 & \texttt{observer.mjs} \\
A window is the multiplier on the torus, with its curvature remainder & any & $9\times 10^{-8}$ & \texttt{observer.mjs} \\
Hard patterns are observer-proof & any & $2\times 10^{-17}$ & \texttt{observer.mjs} \\
Soft patterns reopen a null under squaring, linearly in the ramp & drawing & $3\times 10^{-3}$, $7\times 10^{-3}$, $10^{-2}$ & \texttt{observer.mjs} \\
The octave duty null on paper: the $2{:}1$ station at coarse duty $1/2$ & drawing & $6{,}305\times$ deep & \texttt{dutynull.mjs} \\
The octave duty null in an ear: hard pulses, square-law, cubic, two-stage & sound & $10{,}201\times$, $181\times$, $8{,}564\times$ & \texttt{ear.mjs} \\
A linear ear hears no beat in a sum of trains, at any duty & sound & $4\times 10^{-6}$ of the square-law ear & \texttt{ear.mjs} \\
Softened pulses: a square-law ear keeps the null, a cubic ear reopens it & sound & $132{,}480\times$ against $3.3\times$ & \texttt{ear.mjs} \\
Beats of beats: one stage of degree two or three hears none, a cascade does & sound & $8\times 10^{-5}$, $2\times 10^{-7}$ vs 13\% & \texttt{ear.mjs} \\
A multiplied trio carries the beat of beats at linear order & sound & 50\% of a first-order beat & \texttt{ear.mjs} \\
The golden ratio is a desert in sound & sound & station line $7.5\times$ the golden pair's best & \texttt{ear.mjs} \\
Temporal selection: a shutter keeps the recipes its rates annihilate & a camera & $14{,}668\times$; kept to 0.09\% & \texttt{exposure.mjs} \\
Record beats are convergents; the reduction names the true winner & any & 207 of 207; 99.8\% & \texttt{convergents.mjs} \\
Stations are the convergents with large complete quotients; deserts at $1/\sqrt5$ and $1/2\sqrt2$ & any & 8{,}329 classified; 0.447, 0.354 & \texttt{stations.mjs} \\
Fringe endings count the enclosed winding; the core radius $2As/\pi$ & drawing & 22 probes exact; core within $10^{-4}$ & \texttt{defects.mjs} \\
The fold law: onset radii and the Mach condition & drawing & within 2.8\% & \texttt{foldlaw.mjs} \\
The pooled view is the slide average, exactly & instrument & $7\times 10^{-5}$ over 480 scenes & \texttt{exactsweep.mjs} \\
\bottomrule
\end{tabular}

  \caption{\figlead{The gated predictions, by domain.} ``Any'' means the
  claim is about pictures on a torus and holds for drawings, sounds and
  anything else that counts; ``drawing'', ``sound'' and ``a strobe'' are the
  domains in which it was measured, on drawn strokes, on a model ear, and on
  sampled frames; ``a camera'' is the long exposure; ``instrument''
  certifies the tool. No claim without a number; no number without a
  script.}
  \label{tab:predictions}
\end{table}

\subsection{What is classical}
None of the following is ours, and the paper leans on all of it. That two
sinusoids beat and that a nonlinearity demodulates the beat is
Helmholtz~\cite{Helmholtz1954Sensations}, and every radio detector since.
That the bands of two gratings are the level lines of a difference of
counts is the indicial picture of the moir\'e
literature~\cite{Oster1964, Amidror2009Theory}, which also has a complete
Fourier theory of which orders exist for periodic gratings. The count itself
is old and has many names: the index of the indicial method, the fringe
order or unwrapped phase of interferometry, the phase of an oscillator in
the theory of coupled oscillators, the modulation phase that superspace
crystallography adds as an extra coordinate to describe an incommensurate
crystal~\cite{deWolff1974Modulated}; that an aperiodic family has one is
de Bruijn's index function for multigrids~\cite{deBruijn1981Penrose}, and
the torus of an aperiodic superposition is the hull of the physics of
incommensurate solids~\cite{Bellissard2000Hulls}. What this paper adds
to the count is not the count: it is the factorisation $S=I\circ\Phi$ as
the one object, and what is proved about it below. That averaging along a subgroup is a
conditional expectation is ergodic theory's, and that convergents are the
best approximations and the golden ratio the worst approximable number is
Lagrange's, Hurwitz's and Khinchin's~\cite{Hurwitz1891Annaherung,
Khinchin1964Continued}. Defects with quantized charge are the dislocations
of singular optics~\cite{NyeBerry1974Dislocations}.

\begin{lookup}{what the neighbours already knew}
The names above, in one line each. \emph{Indicial method}: draw each
grating's lines numbered, and the moir\'e bands join the crossings whose
numbers differ by a constant --- \S\ref{sec:coincidence} with a ruler.
\emph{Multigrid index function}: de Bruijn numbered the lines of several
grids to build Penrose tilings; his numbering is a count. \emph{Hull of an
aperiodic solid}: physicists studying two incommensurate lattices describe
where an atom sits in each as a point on a torus. \emph{Ergodic theory}:
the mathematics of averages along a motion, whose basic fact is that the
average along a symmetry keeps exactly what the symmetry leaves fixed.
\emph{Convergents}: the lookup of \S\ref{sec:selection}.
\emph{Dislocations of singular optics}: light waves whose phase winds
round a point, whose fringes fork exactly as \S\ref{sec:counting}'s bands
end.
\end{lookup}

\subsection{What is ours}
The claim of the paper is that these are one subject, and that the one
object $S=I\circ\Phi$ makes them so: every phenomenon above is a property
of the state map, of the picture, or of the way an observer averages one
over the other. The specific contributions are: the sharp conditional
expectation for whole-number slides (Theorem~\ref{thm:sweep}) and the soft
one for windows with its explicit leakage (Theorem~\ref{thm:multiplier});
the observer theorem (Theorem~\ref{thm:universality}) with its corollaries
--- who can see a beat, hard patterns observer-proof, soft ones not ---
and the universal-invariant proposition that makes ``the third pattern''
a quotient (Proposition~\ref{prop:invariant}); selection as best
approximation with the harmonic price, so that stations are places on a
pair whose ratio varies, and the threshold on the complete quotient with
its deserts (\S\ref{sec:selection}); the duty null
(Theorem~\ref{thm:dutynull}) on paper, in a model ear and under a strobe;
the order-four theorem for beats of beats (Theorem~\ref{thm:cascade}); the
trichotomy of counts with the fold law and the defect count
(\S\ref{sec:counting}); and an instrument in which every one of these is a
measurement (\S\ref{sec:instrument}).

\subsection{How sure}
Of the theorems: as sure as a proof of a few lines can make anything; they
are above, and with the lookups in front of them a student can check each
in an afternoon. Of the drawn predictions: as sure as a measurement against
brute force through the shipped renderer, and the table is their record.
Of the predictions in sound and under a strobe: the arithmetic is as sure
as the theorems, and the physics adds assumptions --- that an ear applies
a bending front end before it pools, that a two-stage system exists in it,
that an eye averages consecutive frames --- which are standard but not
ours; the ear of Figure~\ref{fig:ear} is a \emph{model} of an ear, four
lines of code, not a person, and the wheel of \S\ref{sec:selection} is
sampled frames, not a viewer. What the theory does not yet have is a
prediction that a physicist would find surprising and that has been tested
outside the tool on something that is not a simulation. The last
subsection is the list of what would change that, in order of how little it
would cost.

\subsection{Where the depth is, and is not}
The mathematics is elementary: every theorem is a few lines from Fourier
series, the implicit function theorem and the arithmetic of convergents.
That is a strength for a theory meant to be understood and a weakness for
one meant to impress. The depth is in the unification. The same theorem
says why a mistuned octave of square waves does not beat
(Figure~\ref{fig:ear}), why a wheel with $50\%$ sectors cannot show the
doubled still wheel under a strobe (\S\ref{sec:selection}), and why a
blurred grating reopens a null under a squaring observer while a sharp one
does not (Corollary~\ref{cor:soft}); we believe it also says why a
symmetric pair of coupled oscillators has no even locking tongue --- the
frequency ranges in which two oscillators lock at a ratio, a known
symmetry fact of dynamics --- but that one is not measured here and is
listed as a conjecture. A theory that answers questions it was not built
to answer is the standard we are trying to meet; the duty null and the
long exposure's selection of recipes (\S\ref{sec:pooling}) were the first
two such answers, neither sought, both measured the day they were derived.

\subsection{What to test at home, and what it would show}
\begin{enumerate}
\item \emph{The octave null} (Try-it, \S\ref{sec:whosees}). Two pulse
waves at $200$ and $405$ hertz; the five-hertz throb present at $30\%$
duty, absent at $50\%$, present at $70\%$. Cheap; a synthesizer or a web
synth with a pulse-width control. What to expect: at $50\%$ the throb
should be \emph{gone}, not merely weaker --- the model ears put it three to
four orders of magnitude down --- while both tones stay plainly audible.
It would confirm Theorem~\ref{thm:dutynull} in a real ear.
\item \emph{The order of the ear's front end.} The same, with the pulses'
edges softened by a low-pass filter on \emph{each} oscillator, before they
are mixed; a filter after the mixer acts on the sum and is a different
experiment (it is linear, and Theorem~\ref{thm:whosees} says it cannot
change the answer). If the null at $50\%$ survives, the ear's effective
front end is quadratic; if the throb comes back, it is of higher order
(Corollary~\ref{cor:soft}). Cheap, and not in the literature in this form
as far as we can find.
\item \emph{Beats of beats} (Try-it, \S\ref{sec:iterated}). Three sines at
$300$, $330$, $363$ hertz; a three-hertz throb needs a front end of order
four. Listening with the unaided ear is the experiment; building the
two-stage version needs a program that can chain distortion, a low-pass and
distortion again, which any audio workstation can. The psychoacoustics of
``second-order beats'' has data; the theorem says what the data must show.
\item \emph{The desert.} Two sawtooth tones at $200$ and $410$ hertz beat
ten times a second; two at $200$ and $324$ hertz, the golden ratio, sound
rough and do not beat at any rate you can count.
\item \emph{The wagon wheel.} A fan on a phone's video, or a spoked wheel
under a strobe with an adjustable rate. Near the frame rate the wheel turns
slowly and reverses when the spoke rate crosses it. At half the frame rate
a wheel with narrow spokes shows a still wheel with twice the spokes; a
wheel whose spokes are as wide as the gaps between them shows only
flicker (\S\ref{sec:selection}). The frames are simulated and gated; the
eye is the experiment.
\item \emph{The long exposure} (Try-it, \S\ref{sec:pooling}). Two window
screens, one moving, a phone in night mode: the bands survive and the
members do not, and moving both at rates $1{:}2$ keeps a different set of
bands.
\item \emph{The marker} (Try-it, \S\ref{sec:selection}). Two combs drawn
at pitches near $2{:}1$ on paper or transparencies: the octave bands vanish
when the coarse comb's ticks are thickened to half the gap and return at
two thirds. Paper, a pen and a marker; the only test on this list that
needs no electricity.
\end{enumerate}
